\section{Generic metrics and geodesic complexity}
\label{sec:generic}

We establish the general-position statements used above and
give an upper bound for geodesic complexity.
Throughout, $N=2n=\dim(M\times M)$.
All metrics in this section have negative sectional curvature.

We begin by recording the smoothness and variation properties of the distance branches.

Let $g$ be a negatively curved metric on $M$.
The lifted metric on $X=\widetilde M$ is complete,
simply connected and has no conjugate points.
For $(u,v)$ with $u\ne av$, the distance branch
$L_a(u,v)=d_X(u,av)$ is smooth.
Moreover, in a sufficiently small finite-dimensional
smooth family of negatively curved metrics, the branches
$L_a$ vary smoothly in the metric parameters and endpoints:
the unique connecting geodesics vary smoothly by the
geodesic equation and nondegeneracy of the endpoint map.
We use this smoothness only off the diagonal $x=y$,
where different minimizing labels cannot tie.

Two distinct minimizing geodesics from $x$ to $y$
have disjoint interiors. Indeed, if they shared an
interior point, swapping their initial and final
segments would again give a shortest broken path.
A globally minimizing path cannot have a corner,
so uniqueness for the geodesic equation forces
the geodesics to coincide.

\begin{lemma}\label{lem:metric-submersion}
Suppose $a_0,\ldots,a_{k-1}$ represent distinct
minimizing geodesics from $x$ to $y$.
There are smooth variations of the metric,
supported in pairwise disjoint neighborhoods
of interior points of these geodesics, whose
first variations of their lengths form an
invertible diagonal matrix. Consequently,
the universal map
\[
 (g,u,v)\longmapsto
 (L_{a_1}-L_{a_0},\ldots,L_{a_{k-1}}-L_{a_0})
\]
is a submersion at a simultaneous minimum.
\end{lemma}

\begin{proof}
Choose an interior point of each geodesic not
lying on any other minimizing geodesic, and
a nonnegative smooth cutoff $\chi_i$ supported
in a small neighborhood of that point.
Use the symmetric two-tensor variation
$h_i=\chi_i g$. The first variation of the
length of a geodesic with fixed endpoints is
\[
 D_g L_{a_j}(h_i)
 =\frac12\int_{\gamma_j}
    h_i(\dot\gamma_j,\dot\gamma_j)\,ds.
\]
The integral is positive for $i=j$ and
vanishes for $i\ne j$, by disjointness of
the supports. Taking the differences with
the $a_0$ branch gives independent
variations in the target.
\end{proof}

These variations lead to a general-position theorem.

\begin{theorem}\label{thm:genericity}
The set of regular metrics is open and dense
in $\mathcal R_-(M)$.
For a regular metric, the locus of pairs with
exactly $k$ minimizing geodesics is locally a
smooth submanifold of codimension $k-1$.
No such pairs exist when $k>N+1$, and the
locus with $k=N+1$ is finite.
\end{theorem}

\begin{proof}
Fix $g_0\in\mathcal R_-(M)$ and a compact set
$K\subset\widetilde M$ whose deck translates cover
$\widetilde M$. In a sufficiently small $C^\infty$
neighborhood $\mathcal U$ of $g_0$, the lifted metrics
are uniformly bilipschitz to $g_0$, with constant
arbitrarily close to one. Properness of the deck action
and compactness of $K$ therefore give a finite set
$F\subset\Gamma$ such that every minimizing label
at a point of $K\times K$, for every $g\in\mathcal U$,
belongs to $F$. Enlarge $F$, if necessary, to include
all labels whose lengths differ from the minimum by at
most a fixed small allowance. All possible simultaneous
minima are thus drawn from finitely many subsets of $F$.

For a subset $S=\{a_0,\ldots,a_{k-1}\}\subset F$ with
$k\geq2$, consider on the locus where the branches
are smooth the universal difference map
\[
 \mathcal F_S(g,u,v)
 =\bigl(L_{a_1}^g-L_{a_0}^g,\ldots,
        L_{a_{k-1}}^g-L_{a_0}^g\bigr).
\]
At a point where all labels in $S$ are minimizing,
Lemma~\ref{lem:metric-submersion} makes
$\mathcal F_S$ a submersion in the metric parameter
directions. This property is open: after choosing
the finitely many metric variations supplied by the
lemma, their length-variation matrix remains
surjective on a neighborhood of that point.

For the fixed base metric $g_0$, the simultaneous
minimum locus associated with $S$ in $K\times K$ is
compact. Cover it by finitely many such neighborhoods
and collect the metric variations from these
neighborhoods. Their span is a finite-dimensional
space of smooth symmetric two-tensors. Shrink the
metric-parameter ball so that the resulting metrics
remain negatively curved. On a neighborhood of the
simultaneous-minimum locus, the corresponding
finite-parameter map $\mathcal F_S$ is transverse
to zero. On the compact complement of that
neighborhood, simultaneous minimization for $S$
is excluded by a strictly positive length gap;
after a further shrinking of the parameter ball,
it remains excluded. Thus finite-dimensional
parametric transversality produces arbitrarily
small perturbations for which the fixed-label
difference map is transverse to zero at every
simultaneous minimum for $S$.

Repeat the construction for the finitely many
subsets $S\subset F$ and take the span of all
the resulting metric variations. Transversality
for one configuration is stable under adjoining
the other parameters, because the earlier
surjective parameter directions remain available.
The parametric transversality theorem, applied
to all configurations simultaneously, gives
metrics arbitrarily close to $g_0$ satisfying
the regularity condition. This proves density.

For openness, suppose $g_j\to g$ and that every
$g_j$ is nonregular, whereas $g$ is regular.
After putting the corresponding lifts of endpoint
pairs in $K\times K$, the finite-label reduction
allows us to pass to a subsequence with a fixed
set $S$ of at least two minimizing labels witnessing
failure of full rank. Compactness gives a limit
$(u,v)$ at which all labels in $S$ still minimize
for $g$. The non-full-rank condition is closed
under convergence of the differentials, contradicting
regularity of $g$. Hence regular metrics form an
open subset.

At a regular pair with exactly $k$ minimizing
labels, the implicit function theorem cuts out
the equal-length locus with codimension $k-1$.
All other labels have a positive length gap in a
small neighborhood, so this is also the local
exact-multiplicity stratum. The inequality
$k-1\leq N$ follows from full rank. The stratum
of multiplicity $N+1$ has dimension zero.
It is closed, since the existence of at least
$N+1$ minimizing labels is a closed condition
under the finite-label reduction; compactness
of $M\times M$ then makes it finite.
\end{proof}

We next establish a selection result for minimizing branches on the multiplicity strata.

We recall that $\GC(M,g)$ is the minimum
number of locally compact pieces partitioning
$B=M\times M$, each with a continuous choice
of minimizing geodesic, using the unreduced
convention of \cite{RecioMitter2021}.
The following selection argument is useful
more generally for a finite-valued covering
with clean multiplicity strata.

\begin{lemma}\label{lem:selection}
Let $B$ be a second-countable smooth $N$-manifold
(without boundary), let $q:E\to B$ be a smooth
covering, and let $Z\subset E$ be closed. Assume
$q|_Z$ is proper and surjective, with nonempty
finite fibers. For $k\geq1$ put
\[
 B_k=\{b\in B:\#(Z\cap q^{-1}(b))=k\}.
\]
Assume every nonempty $B_k$ is a smooth
embedded submanifold of $B$ (not necessarily
closed), of dimension at most $N-k+1$.
Then $B$ admits a partition into at most
$N+1$ pairwise disjoint locally closed subsets
on each of which $q|_Z$ admits a continuous
section. In particular, the conclusion does
not assume that the finite-valued map
$q|_Z$ is a covering.
\end{lemma}

\begin{proof}
Because $q|_Z$ is proper and $Z$ is closed,
the fiber cardinality $k(b)$ is upper
semicontinuous. Indeed, given $b\in B$, choose
an evenly covered neighborhood and disjoint
local sheets through the finitely many points
of $Z_b$. Properness excludes any other
$Z$-points over a sufficiently small
neighborhood of $b$. Thus nearby fibers
are subsets of these finitely many sheets.
Consequently $B_k$ is locally closed;
the smooth embedded structure assumed in
the statement is compatible with the
subspace topology. Moreover $k\leq N+1$
whenever $B_k$ is nonempty.

Write $d_k=\dim B_k\leq N-k+1$. We first
choose $N+1$ smooth functions
$f_0,\ldots,f_N:E\to\R$ with a simultaneous
general-position property on all the strata.

Fix $k$ and a point of $B_k$. Since $q$ is a
covering and $Z\to B$ is proper with finite
fibers, after shrinking a neighborhood of the
point we can enumerate the $k$ active sheets
by smooth local sections of $q$. Restricting
these sections to $B_k$, we obtain smooth
maps $z_1,\ldots,z_k$ with pairwise disjoint
images in $E$. For a tuple
$f=(f_0,\ldots,f_N)$, consider the evaluation
map
\[
 F_f:B_k\longrightarrow (\R^k)^{N+1},
 \qquad
 F_f(b)=\bigl((f_j(z_i(b)))_{i=1}^{k}
                  \bigr)_{j=0}^{N}.
\]
For each collection
$\mathcal P=(\mathcal P_0,\ldots,\mathcal P_N)$
of partitions of $\{1,\ldots,k\}$, let
$D_{\mathcal P}\subset(\R^k)^{N+1}$ be
the linear subspace imposing equality of
coordinates inside each block of
$\mathcal P_j$, for each $j$. Its codimension
is
\[
 \operatorname{codim}D_{\mathcal P}
 =\sum_{j=0}^N(k-|\mathcal P_j|).
\]
The parametrized evaluation map
$(f,b)\mapsto F_f(b)$ is a submersion in the
function variables: values at the finitely
many distinct points $z_i(b)$ can be
perturbed independently with disjointly
supported smooth bump functions on $E$.
Parametric transversality therefore gives
a residual set of tuples $f$ for which
$F_f$ is transverse to every
$D_{\mathcal P}$ on the chosen chart.

There are only finitely many partitions for
each $k$, and countably many sheets and
charts suffice to cover the second-countable
strata $B_k$. Using compact exhaustions
and the Baire property of the space of
smooth functions on $E$, choose one tuple
$f$ transverse to all these diagonals
simultaneously. When
$\operatorname{codim}D_{\mathcal P}>d_k$,
transversality implies that
$F_f^{-1}(D_{\mathcal P})$ is empty.
Consequently, at every $b\in B_k$, if
$\mathcal P_j(b)$ is the partition of the
active $k$ sheets according to equality of
their $f_j$-values, then
\begin{equation}\label{eq:selection-inequality}
 \sum_{j=0}^N
 \bigl(k-|\mathcal P_j(b)|\bigr)
 \leq d_k\leq N-k+1.
\end{equation}

Define
\[
 c(b)=k(b)-1+
 \sum_{j=0}^N
 \bigl(k(b)-|\mathcal P_j(b)|\bigr).
\]
Thus $0\leq c(b)\leq N$.
The vectors
\[
 (f_0(z),\ldots,f_N(z)),\qquad z\in Z_b,
\]
are pairwise distinct. Otherwise two
active sheets would have equal values for
all $N+1$ functions, forcing each term
$k-|\mathcal P_j|$ to be at least one,
contrary to \eqref{eq:selection-inequality}.
Hence there is a unique lexicographically
smallest vector, determining a selected
sheet over each $b$.

We verify that the selection is continuous
on the sets $A_j=\{b:c(b)=j\}$.
Fix $b$ and enumerate all $k(b)$ active
sheets in a covering neighborhood $U$.
Properness of $Z\to B$ and closedness of
$Z\subset E$ allow $U$ to be chosen so
that every active sheet at $b'\in U$
is among these enumerated sheets.
Furthermore, pairs with unequal
$f_j$-values at $b$ remain unequal after
shrinking $U$. Therefore, as $b'\to b$,
the number of active sheets cannot
increase, and each equality partition can
only split or lose elements. More precisely,
for every $j$ the nonnegative integer
$k(b')-|\mathcal P_j(b')|$ cannot exceed
its value at $b$, and $k(b')\leq k(b)$.
Hence $c$ is upper semicontinuous.

Since $c$ is integer-valued with range
$\{0,\ldots,N\}$, each
\[
 A_j=\{c\geq j\}\cap\{c<j+1\}
\]
is locally closed and thus locally compact.
Fix $b$ and shrink $U$ as above. For
$b'\in A_{c(b)}\cap U$, all nonnegative
contributions
\[
 k(b')-1,\qquad
 k(b')-|\mathcal P_j(b')|
 \quad(0\leq j\leq N)
\]
are bounded above by the corresponding
ones at $b$. Their sums are equal, so
each contribution agrees. Hence all
$k(b)$ active sheets persist and the
equality partitions have not split.
Every pair of sheet values that was
unequal at $b$ remains unequal with the
same sign in a sufficiently small
neighborhood. Thus the lexicographic
order between any two distinct evaluation
vectors is constant along
$A_{c(b)}\cap U$. Its unique minimum
therefore belongs to the same covering
sheet as at $b$, proving continuity
of the selected section on $A_{c(b)}$.
The disjoint locally closed sets
$A_0,\ldots,A_N$ therefore give the
required partition and continuous
sections.
\end{proof}

\begin{theorem}\label{thm:generic-GC}
For every regular negatively curved metric
on $M$ one has
\[
 \GC(M,g)\leq 2n+1.
\]
If $M$ is closed and oriented, then
\[
 \GC(M,g)=\TC(M)=2n+1
\]
and
\[
 \GC(M,g)
 =1+\max_{x\ne y}
   \operatorname{rank}\MH_{2,d_M(x,y)}(M;x,y).
\]
\end{theorem}

\begin{proof}
Let $E\to B$ be the covering whose fiber
over $(x,y)$ consists of fixed-endpoint
homotopy classes of paths, equivalently
the labels of lifts. Let $Z$ be the subset
corresponding to minimizing geodesics.
The map $Z\to B$ is closed and proper with
finite fibers, by compactness of $M$ and
properness of the universal covering action.
The exact multiplicity strata are smooth
of dimension $N-k+1$, by
Theorem~\ref{thm:genericity}.
Apply Lemma~\ref{lem:selection}. A continuous
selection of a minimizing label gives a
continuous minimizing geodesic, because
the corresponding lifted geodesic depends
smoothly on its endpoints. Thus
$\GC(M,g)\leq N+1$.

The inequality $\TC(M)\leq\GC(M,g)$ is
standard \cite{RecioMitter2021}. The
known calculation of topological complexity
for torsion-free hyperbolic groups
\cite{Dranishnikov2020}, translated into
our unreduced convention, gives
$\TC(M)=N+1$. Hence equality holds.

For the final assertion, we do not use the
quantitative counting theorem. Horja's
multiplicity theorem \cite{Horja1997}
shows that a pair realizing a local maximum
of the distance function on a closed
negatively curved $n$-manifold is joined
by at least $2n+1=N+1$ minimizing
geodesics. Such a pair exists by compactness.
For a regular metric, Theorem~\ref{thm:genericity}
gives the reverse multiplicity bound $N+1$.
Hence the maximal number of minimizing
geodesics is exactly $N+1$. By
Proposition~\ref{prop:MH-minimal}, the
largest minimal-length degree-two magnitude
homology rank is $N$.

In particular, the geodesic-complexity
calculation and this rank identity are
independent of the quantitative seminorm
and localization estimates in
Sections~\ref{sec:seminorm} and
\ref{sec:localization}.
\end{proof}
